% !TEX root = ../macro.tex
\clearpage
\part{Topics}\label{part:topics}

The topics combine mechanisms from the earlier parts. Each begins with dated
facts and measurement choices, compares competing explanations, and asks what
evidence would distinguish them. Policy conclusions remain conditional on
that evidence.

\clearpage
\plannedsection{The Dollar in the International Monetary System}\label{sec:dollar}

\subsection{History and measures of international currency use}
Separate reserve holdings, invoicing, payments, funding, and exchange-rate
anchors. Define what each measure can establish about dollar dominance.

\subsection{Network effects, invoicing, and safe assets}
Explain the mechanisms that can sustain an international currency's use.
Connect dollar pricing and demand for safe assets to the open economy models.

\subsection{Benefits, obligations, and the international balance sheet}
Distinguish seigniorage, financing advantages, and valuation effects.
Examine the obligations and constraints associated with supplying safe and
liquid international liabilities.

\subsection{Global funding and policy spillovers}
Trace dollar borrowing, funding stress, and the international transmission of
US policy. Examine liquidity facilities and the role of financial sanctions.

\subsection{Persistence, alternatives, and testable changes}
Compare explanations of persistence with possible sources of change.
Assess claims about de-dollarization separately across the measures defined
at the start of the chapter.

\clearpage
\plannedsection{China's Growth and Adjustment after the Pandemic}\label{sec:china}

\subsection{The baseline, timeline, and measurement}
Separate the pre-pandemic baseline, pandemic years, and reopening period.
Distinguish output levels from growth rates and real activity from nominal
income. An observation after the pandemic does not identify a pandemic effect.

\subsection{Demand, consumption, and employment}
Compare cyclical explanations using household income, consumption, prices,
and employment. State what each explanation predicts before assessing it.

\subsection{Property, balance sheets, and local public finance}
Map the links between housing, household and firm balance sheets, banks,
and local government finances. Compare expenditure adjustment and debt
resolution mechanisms.

\subsection{Population, productivity, and resource allocation}
Use growth accounting and the growth models to distinguish demographic,
investment, and productivity explanations. Examine the allocation of resources
across sectors and firms.

\subsection{External demand and policy alternatives}
Connect domestic adjustment to trade and international conditions.
Compare policy alternatives through their mechanisms, constraints, distributional
effects, and testable implications.
